\documentclass{handout}
\usepackage{handoutSetup}\usepackage{handoutShortcuts}  
\begin{document}
\handoutHeader

\begin{verbatimwrite}{\jobname.title}
The Consumption Capital Asset Pricing Model (C-CAPM)
\end{verbatimwrite}

\handoutNameMake

Consider a representative agent solving the joint consumption and 
portfolio allocation problem:
\begin{equation}\begin{gathered}\begin{aligned}
        \vFunc(\mRat_{t}) & =  \max \uFunc(c_{t}) + \Ex_{t}\left[\sum_{n=1}^{\infty} \Discount^{n} \uFunc(c_{t+n}) \right]
\\  & \text{s.t.}  \nonumber
\\      \mRat_{t+1} & =  (\mRat_{t}-c_{t})\Rport_{t+1} + y_{t+1}  
\\      \Rport_{t+1} & =  \sum_{i=1}^{m}\omega_{t,i}\Risky_{t+1,i}+\left(1-\sum_{i=1}^{m} \omega_{t,i}\right)\Rfree
\end{aligned}\end{gathered}\end{equation}
where $\Rfree$ denotes the return on a perfectly riskless asset and 
$\Risky_{t+1,i}$ denotes the return on asset $i$ between periods $t$ and 
$t+1$, $\omega_{t,i}$ is the share of end-of-period savings invested in 
asset $i$, and $\Rport_{t+1}$ is the portfolio-weighted rate of 
return, and $y_{t+1}$ is noncapital income in period $t+1$.

As usual, the objective can be rewritten in recursive form:
\begin{equation}\begin{gathered}\begin{aligned}
        \vFunc(\mRat_{t}) & =  \max_{\{c_{t},\omega_{1,t},\omega_{2,t},\ldots\}} \uFunc(c_{t}) +\beta \Ex_{t}\left[\vFunc(\Rport_{t+1}(\mRat_{t}-c_{t})+{y}_{t+1})\right]
\end{aligned}\end{gathered}\end{equation}

The first order condition with respect to $c_{t}$ is
\begin{equation}\begin{gathered}\begin{aligned}
        \uFunc^{\prime}(c_{t}) & =  \beta \Ex_{t}[\Rport_{t+1}\vFunc^{\prime}({m}_{t+1})] \label{eq:cfoc}.
\end{aligned}\end{gathered}\end{equation}
and the FOC with respect to $\omega_{t,i}$ is
\begin{equation}\begin{gathered}\begin{aligned}
        \Ex_{t}[(\Risky_{t+1,i}-\Rfree)\vFunc^{\prime}({m}_{t+1})] & =  0 \label{eq:gamfoc}
\end{aligned}\end{gathered}\end{equation}

But the usual logic of the Envelope theorem tells us that 
\begin{equation}\begin{gathered}\begin{aligned}
        \uFunc^{\prime}(c_{t+1}) & =  \vFunc^{\prime}(\mRat_{t+1}), \label{eq:envelope}
\end{aligned}\end{gathered}\end{equation}
so, substituting (\ref{eq:envelope}) into (\ref{eq:cfoc}) and (\ref{eq:gamfoc}) we have
\begin{equation}\begin{gathered}\begin{aligned}
        \uFunc^{\prime}(c_{t}) & =  \Ex_{t}\left[\beta \Rport_{t+1} \uFunc^{\prime}({c}_{t+1})\right] %\label{eq:ceuler}
\\      \Ex_{t}[(\Risky_{t+1,i}-\Rfree)\uFunc^{\prime}({c}_{t+1})] & =  0 \label{eq:gameuler}
.
\end{aligned}\end{gathered}\end{equation}

Now assume CRRA utility, $\uFunc(c) \equiv c^{1-\rho}/(1-\rho)$ and divide both 
sides of (\ref{eq:gameuler}) by $c_{t}^{-\rho}$ to get
\begin{equation}\begin{gathered}\begin{aligned}
        \Ex_{t}[(c_{t+1}/c_{t})^{-\rho}(\Risky_{t+1,i}-\Rfree)] & =  0 \label{eq:gameulernew}
\end{aligned}\end{gathered}\end{equation}


We can now follow the same steps as in the `Equity Premium Puzzle' handout to
obtain the relation that {\it for every asset $i$} 

\begin{equation}\begin{gathered}\begin{aligned}
        \Ex_{t}[\Risky_{t+1,i}]-\Rfree & \approx  \frac{\rho \text{cov}_{t}(\Delta \log {c}_{t+1},\Risky_{t+1,i})}{1-\rho \Ex_{t}[ \Delta \log {c}_{t+1}]} 
\\                             & \approx  \rho \text{cov}_{t}(\Delta \log {c}_{t+1},\Risky_{t+1,i}) \label{eq:eqprem}
\end{aligned}\end{gathered}\end{equation}


What does this imply about asset pricing?  

Consider an asset for which the return covaries positively with 
consumption, $\mbox{cov}(\Risky_{t+1,i},\Delta \log c_{t+1}) > 0$.  For 
such an asset, the {\it marginal utility} will negatively covary with 
the return.  Thus the expected return must be higher for an asset 
that `does well' when consumption is high.  But for a given average 
stream of dividends or payouts, if the average return is high, the 
average price must be low.  Thus this indicates that prices should be 
low for assets whose payoffs are procyclical, and high for assets 
whose payoffs are countercylical.





\end{document}

 
\input handoutBibMake
